\section{Complete optimal-prefix reconstruction at smaller order ten}
\label{sec:order-ten}

We now close the finite range left by \cref{thm:eventual}.  The argument
is an exact certificate over one prime field.  Its force comes from the
analytic reduction already proved in the source: computation is used
only after the problem has become a finite list of explicitly specified
rational determinants.

\subsection{A good prime and exact Bernoulli reduction}

Fix
\begin{equation}
 P=1{,}000{,}003.
 \label{eq:certificate-prime}
\end{equation}
Trial division by the odd integers through $1000$ proves that $P$ is
prime; this finite check is included in both verifiers.  The Bernoulli
numbers are reduced modulo $P$ from their defining recurrence, not from
a table.  With $B_1=-1/2$ and $m\ge1$,
\begin{equation}
 B_{2m}=\frac{\dfrac{2m-1}{2}
       -\displaystyle\sum_{j=1}^{m-1}\binom{2m+1}{2j}B_{2j}}
       {2m+1}.
 \label{eq:even-bernoulli-recurrence}
\end{equation}
This follows immediately from
$\sum_{j=0}^{2m}\binom{2m+1}{j}B_j=0$ and the vanishing of odd
Bernoulli numbers beyond $B_1$.

For the retained matrices in the range $10\le b\le1919$, the largest
factorial argument occurring in a matrix entry is $3838$, the largest
denominator appearing in the Bernoulli recurrence is $3839$, and the
largest Bernoulli number used by an entry is $B_{3836}$.  The scalar
multiplier in \eqref{eq:exact-multiplier} contains factorials only up to
$7674$.  All these integers are strictly below $P$.  It follows
inductively from \eqref{eq:even-bernoulli-recurrence} that every rational
Bernoulli number occurring in the certificate has denominator prime to
$P$.  Every factorial and every remaining explicit linear denominator
is likewise a unit in $\F_P$.  Thus all source formulas admit an
unambiguous reduction modulo $P$.

\begin{lemma}[Denominator-safe reduction]
\label{lem:denominator-safe}
For $6\le b\le9$, reduction of the source matrix $M_{5,b}$ modulo $P$
is defined entrywise.  For $10\le b\le1919$, reduction of the retained
matrix $\A^{[5]}$ and of the multiplier in
\eqref{eq:exact-multiplier} is defined, and the multiplier is nonzero in
$\F_P$.  Consequently,
\[
 \det\A^{[5]}\not\equiv0\pmod P
 \quad\Longrightarrow\quad
 \det M_{5,b}\ne0\quad\text{in }\Q.
\]
The analogous implication with the direct determinant holds for
$6\le b\le9$.
\end{lemma}

\begin{proof}
For the four direct cases the largest factorial argument is $25$.  In
the retained range, inspection of \eqref{eq:retained-indices}--
\eqref{eq:retained-A} gives the bounds stated above.  Every factor in
\eqref{eq:exact-multiplier} is a product or quotient of nonzero
integers less than $P$, so it is nonzero in $\F_P$.  Reduction modulo
$P$ is therefore a ring homomorphism on the localization of $\Z$
generated by all denominators in question.  A zero rational determinant
would reduce to zero; the contrapositive proves the claim.
\end{proof}

\subsection{The exhaustive residue vector}

Define the proof residue $r_b\in\F_P$ by
\begin{equation}
 r_b=
 \begin{cases}
  \det M_{5,b}\pmod P,&6\le b\le9,\\[2mm]
  \det\A^{[5]}(b;2B_{2b-2},\ldots,2B_{2b-18})\pmod P,
      &10\le b\le1919.
 \end{cases}
 \label{eq:proof-residue}
\end{equation}
The split at $b=10=2a$ is important: below it, the generic retained
formula has endpoint terms not covered by its derivation, so the source
matrix is used directly.

\begin{theorem}[One-prime finite completion]
\label{thm:finite-order-ten}
For every integer $b$ with $6\le b\le1919$,
\begin{equation}
 r_b\ne0\quad\text{in }\F_{1{,}000{,}003}.
 \label{eq:all-residues-nonzero}
\end{equation}
Hence $\det M_{5,b}\ne0$ for every such $b$.
\end{theorem}

\begin{proof}
The exact program \texttt{generate\_order10\_certificate.py} performs
the following finite calculation.

\begin{enumerate}
\item It constructs factorials and inverse factorials through $3839$
in $\F_P$ and generates all required Bernoulli residues from
\eqref{eq:even-bernoulli-recurrence}.
\item For $6\le b\le9$ it builds \eqref{eq:source-M}; for
$10\le b\le1919$ it builds the five-square matrix
\eqref{eq:retained-A} at \eqref{eq:U-specialization}.
\item It computes each determinant by exact Gaussian elimination in
$\F_P$ and records all $1914$ residues.
\end{enumerate}

Every recorded residue lies in $\{1,\ldots,P-1\}$.  The complete vector
is contained in \nolinkurl{verification/order10_modular_certificate.json}.  Its
canonical comma-separated encoding has SHA-256 digest
\begin{equation}
\begin{split}
&\texttt{3977e54f5b43db43e8f080c32f2e5b22}\\[-1mm]
&\hspace{21mm}\texttt{beac54412cca32ff5d03d84975111836}.
\end{split}
\label{eq:residue-digest}
\end{equation}
For orientation, five entries are
\begin{center}
\begin{tabular}{crrrrr}
\toprule
$b$&6&8&10&962&1919\\
\midrule
$r_b$&484126&572800&668726&989617&667530\\
\bottomrule
\end{tabular}
\end{center}
The value at $b=962$ lies in the exceptional $2$-adic progression left
open by the preceding order-ten report; it is nonzero here.

A second verifier replays the entire calculation without importing the
first implementation.  It uses the full recurrence
$\sum_{k=0}^n\binom{n+1}{k}B_k=0$ for every Bernoulli number and a
row-by-row subset expansion for determinants, rather than the even-only
recurrence and Gaussian elimination.  It reproduces all $1914$ residues
and the digest \eqref{eq:residue-digest}.  The digest is an integrity
check, not a substitute for the supplied vector and replay programs.
Finally, \cref{lem:denominator-safe} converts each modular nonvanishing
statement into rational nonvanishing.
\end{proof}

The implementation also checks the complete multiplier
\eqref{eq:exact-multiplier} at $b=10,11,12,13,14$ by comparing the
direct determinant with the retained determinant times every scalar
factor.  These checks defend against an implementation normalization
error; the mathematical justification for all $b\ge10$ is the inherited
exact multiplier lemma.

\subsection{The completed fixed-order theorem}

\begin{theorem}[Complete smaller-order-ten reconstruction]
\label{thm:complete-order-ten}
For every integer $b\ge6$,
\begin{equation}
 \boxed{\det M_{5,b}\ne0.}
 \label{eq:complete-order-ten}
\end{equation}
Equivalently, for every even $q>10$, the first $q$ centered moments for
the harmonic-tail pair $(p,q)=(10,q)$ have full formal rank $q$ on the
specified coordinate list
\[
 C_0,\ldots,C_{q/2-1},\quad
 \zeta(0),\zeta(2),\ldots,\zeta(q-2).
\]
\end{theorem}

\begin{proof}
The finite interval $6\le b\le1919$ is
\cref{thm:finite-order-ten}.  The half-line $b\ge1920$ is
\eqref{eq:eventual-order-ten}.  The rank conclusion follows from
\cref{prop:rank-criterion}.
\end{proof}

\begin{remark}[What the theorem does and does not say]
The result completes one infinite fixed-order family and resolves the
first explicit task in the research agenda of the quartic reconstruction
report.  It does not prove the unrestricted conjecture for every
$a<b$, and it does not assert arithmetic independence of even zeta
values.  The proof is conventional exact mathematics with a finite
computer-verified lemma; it has not been formalized in a proof assistant
or independently refereed.
\end{remark}
